\section{Introduction}

\subsection{Background and Motivation}

Solubility, whether in water or in organic media, is a physicochemical property relevant to medicinal chemistry, pharmaceutics, and chemical-process engineering. During early discovery and developability assessment, more than 40\% of candidate active pharmaceutical ingredients (APIs) show restricted absorption and poor bioavailability because their solid--liquid solubility is too low, which directly curtails their viability for clinical development\cite{delaney2004,solubility2024}. Downstream, in recrystallisation purification, solvent-system selection, and liquid-formulation design, the continuous solubility behaviour of a target molecule across many industrial solvents and operating temperatures governs supersaturation control, product yield, polymorphic purity, and the attainable processing window\cite{vermeire2022,fastsolv2025}.

Faced with the vast condition space formed by a flood of novel drug scaffolds, dozens of common organic solvents, and a continuous temperature gradient, conventional trial-and-error measurements consume scarce starting materials and synthesis time, and they cannot offer full coverage once the space explodes in size. Building an intelligent model that estimates conditional solubility from the drug structure, the solvent structure, and the operating temperature, using cheminformatics representations and machine learning, would speed up high-throughput virtual screening and improve process decisions\cite{boobier2020,ye2021}.

\subsection{Key Challenges}

Although machine learning has accumulated a rich body of work on aqueous or single-temperature solubility estimation\cite{delaney2004,solubility2024}, multi-solvent, multi-temperature prediction in realistic pharmaceutical settings remains difficult for several reasons.

Cross-structure generalisation to unseen parents confronts a leakage hazard inherent to chemical data. A predictive system has to extrapolate to novel drug scaffolds that never appeared in the training set. Yet such chemical datasets routinely contain dense repeated measurements of the same drug in many solvents and at many temperatures. Row-wise random splits in cross-validation then allow records of one parent to seep into both the training and the test folds, which inflates the metrics and gives a misleading impression of accuracy\cite{leakage2023,solubility2024}. A held-out scheme that physically isolates parent scaffolds is needed to measure generalisation honestly\cite{varma2006,cawley2010}.

At the same time, the long-tailed distribution of observational data aggravates the learning bias. A handful of classic drugs (paracetamol among them) contribute dozens of measurements under different conditions, whereas numerous rare molecules appear only in a few sporadic records. An unweighted mean-squared-error objective is inevitably dominated by the frequent parents, so the learner overfits the common scaffolds and sacrifices accuracy on the rare ones.

Further complicating matters are the nonlinear physical coupling of temperature and the confounding embedded in observational data\cite{pearl2009}. The response of solid--liquid solubility to temperature is governed tightly by solute--solvent interaction thermodynamics. In the non-random observational data used here, the testing temperature may correlate with how insoluble a compound is (poorly soluble species tend to be measured at higher temperatures), and this hidden confounding can even produce a spurious negative correlation in the pooled data that contradicts endothermic dissolution, a Simpson's paradox. Separating a constant structural baseline from a continuous temperature response is the central problem this study addresses\cite{vermeire2022,sensitivity2020}.

\subsection{Objectives and Technical Routes}

To tackle these difficulties, this study targets extrapolative conditional-solubility prediction and, under a nested cross-validation framework that groups strictly by parent\cite{varma2006}, constructs and compares two technical routes that blend physical insight with data-driven methods.

For physically interpretable temperature responses, we introduce SOLAR-H (Structured Orthogonal Linear \& Additive Regression -- Hybrid), a physics-informed statistical decomposition. Anchored in the classical van 't Hoff thermodynamics\cite{vermeire2022}, it splits the conditional log-solubility explicitly into a nonlinear structural baseline and a local linear temperature-response term. Within-combination de-meaning (fixed effects) removes the temperature-invariant structural constants, so that the local temperature slope is estimated without bias from purely within-group temperature variation; a nested, adaptive temperature-gain shrinkage then curbs overfitting when extrapolating across parents.

For accurate fitting in high-dimensional chemical space, we build ExtraTrees--TabPFN, a model rooted in feature augmentation and heterogeneous blending. Following the like-dissolves-like principle, it assembles an 869-dimensional feature space that spans 217 physicochemical descriptors plus explicit solute--solvent interactions (products and differences)\cite{boobier2020}; injects parent-equal sample weights ($q=1$) into the tree-splitting criterion to counter long-tailed sampling bias\cite{extratrees2006}; and incorporates the recent pretrained tabular foundation model TabPFN v2\cite{tabpfn2025}, whose optimal non-negative convex blending weight is selected on inner out-of-fold predictions so that orthogonal tree splits and smooth attention manifolds complement each other.

\subsection{Contributions}

This work spans theoretical modelling, algorithmic refinement, mechanistic analysis, and system implementation, organised along the following axes.

On evaluation and optimisation, we standardise 161 chemical parents and design a nested five-outer/three-inner parent-grouped cross-validation (Nested GroupKFold) that prevents data leakage\cite{leakage2023,varma2006}, and we define a parent-weighted empirical risk together with a sample reweighting function, removing the systematic bias of long-tailed sampling from both the optimisation and the evaluation sides.

On algorithms, we build and ablate both SOLAR-H and ExtraTrees--TabPFN, confirming that the explicit intermolecular interactions (products and absolute differences), the parent weighting, and the complementary TabPFN branch each contribute measurably to lower error. In a strict test over 128 unseen parents (1507 development records), the blended model attains a parent-weighted MAE of 0.6668 and an $R^2$ of 0.5234, cutting the error of a conventional tree baseline by 8.54\% (paired 95\% interval $[-0.1219, -0.0082]$\cite{extratrees2006,tabpfn2025}).

On interpretability and causality, TreeSHAP quantifies the feature-group contributions and shows that explicit molecular interactions (attribution 0.432) and drug descriptors (0.466) dominate\cite{treeshap2020}; a within-combination fixed-effects analysis uncovers the genuine upward physical trend of solubility with temperature (within-group slope $+0.01538$), resolving the Simpson's paradox in the pooled data, and a Cinelli--Hazlett partial-$R^2$ sensitivity analysis shows the result is robust to unmeasured confounding up to a strength of 0.54\cite{pearl2009,sensitivity2020}).

On implementation, a full-featured interactive prediction platform built on FastAPI and a modern front end offers single-point queries, side-by-side model comparison, continuous temperature-response scanning, coverage and applicability-domain warnings, and a high-throughput CSV batch engine; automated tests confirm that the online service and the offline models agree to zero numerical error, with sub-second response under hardware acceleration.

\subsection{Organisation of the Report}

Section 2 reviews prior work on solubility prediction, molecular representation, tabular learners, and evaluation. Section 3 formalises the problem, the data, parent standardisation, and the nested evaluation. Section 4 derives and ablates the physics-based SOLAR-H model. Section 5 builds the ExtraTrees--TabPFN model around high-dimensional explicit interactions and heterogeneous blending. Section 6 compares the two models on held-out parents, with stratified errors and robustness checks. Section 7 dissects the decision mechanisms and the true temperature effect through TreeSHAP attribution and fixed-effects causal analysis. Section 8 describes the interactive prediction system. Section 9 summarises the contributions, the limitations, and the roadmap.
